\clearpage
\fancyhead[L]{\small EFFEMERIDI ASTRONOMICHE.}
\SourceLeaf{561}
\begin{center}\small Göttingische gelehrte Anzeigen. No. 95. Pages 945..952. 14 June 1813.\\[3pt]\rule{.72\linewidth}{.3pt}\end{center}

\emph{Effemeridi astronomiche di Milano per l'anno bisestile} 1812 \emph{calcolate da} \textsc{Francesco Carlini e Carlo Brioschi}. \emph{Con appendice.} 1811. Della reale stamperia. Milano. The calendar has 124 pages, and the supplements just as many, in small quarto.

\emph{Effemeridi astronomiche di Milano per l'anno} 1813, and so forth. 1812. The calendar has 95 pages and the supplements 136 pages, in small quarto.

The Milan astronomical ephemerides are little known here, though their exemplary arrangement and substantial supplements rival those of every similar work currently appearing. It is a pity that they always appear only a short time before the beginning of the year for which they are intended and therefore usually reach us only after that year has begun. If only the editors would follow Mr. \textsc{Bode}'s example in this respect! He always brings out his yearbooks more than two years ahead of time.

The calendar is arranged as follows. Besides a key to the abbreviations and some notes on chronology and church festivals, the preliminary material gives a general notice of eclipses, the apparent obliquity of the ecliptic, and the nutation in longitude of the equinoctial points, both of the latter including solar nutation. The months follow, with 6 pages devoted to each. The first gives the phases of the Moon, conjunctions of the Moon with fixed stars, actual occultations of those stars by the Moon, close conjunctions of planets with fixed stars, and other noteworthy moments in the motions of the planets and the Sun. It also gives the eclipses of Jupiter's satellites in mean time. The second page shows how far each day is from the beginning of the year, the days of the week, mean time and sidereal time at apparent noon, sidereal time at mean noon, and sunrise and sunset. The third contains the Sun's longitude, calculated with the greatest care to tenths of a second, its right ascension in arc, its declination, and the logarithms of its distance from Earth to six decimal places. The logarithms are given for every single day, a very considerable advantage peculiar to these ephemerides. The fourth and fifth pages are devoted to the Moon and give its longitude, latitude, horizontal parallax, and diameter for both noon and midnight, the times of its culmination, rising, and setting (only in minutes of time), and its declination at meridian transit in arcminutes. Finally, the sixth page shows the configurations of Jupiter's satellites for each day of the month. The Sun's semidiameter, the time it takes to cross the meridian, and the longitude of the Moon's node are then given at 6-day intervals throughout the year. Next come the motions of each of the ten planets. We find it very convenient to have these together in a single overview for the whole year instead of divided up by month. Mercury, Venus, and Mars each take up two pages, as their entries are given every 6 days throughout the year. Ceres, Pallas, Juno, Vesta, Jupiter, Saturn, and Uranus take up only half as much space, since the positions of the first four are given only for the periods when they are visible, and those of the others only every 12 days. For each planet the longitude, latitude, right ascension (in minutes of time), declination, rising, culmination, and setting are given to the minute. If we may make one further request, it would be that the right ascensions be given to the nearest arcminute (especially important for the new planets) and that a separate column be devoted to the planets' distances from Earth. These would be useful especially for calculating aberration, and also for \SourceLeaf{562}parallax, apparent diameters, and other purposes. To make room, the rising or setting of the planets could be omitted according as they culminate in the evening or early morning. The longitude column could also be made somewhat narrower by printing longitudes in degrees and minutes alone instead of in signs, degrees, and minutes. We also note that a calculation error makes all the right ascensions of Pallas in 1813 wrong, as well as the culmination times that depend on them, and that for the former the complements to 42 hours must be taken. --- Finally, occultations of fixed stars are given from the calculations of the Florentine astronomers, just as they are also printed every year in the Monatl. Correspondenz. Besides these regular items, the calendar for 1812 contains a list of all stars above the fifth magnitude visible from Milan, reduced to 181c according to \textsc{Piazzi} (596 in all), and \textsc{Gauss}'s tables for aberration and nutation, along with a small table for solar nutation.

Among the \emph{supplements} to the volume for 1812, the first place goes to \textsc{Barnabas Oriani}'s observations of zenith distances of the Sun and fixed stars on the meridian, made with a new repeating circle. They are preceded by an instructive description of the incomparable instrument, a three-foot \textsc{Reichenbach} circle with an upright column (price 3000 guilders). The observations were made with this instrument over two years and printed in the volumes for 1812 and 1813. They offer a wealth of experience for stellar declinations, the theory of the Sun's motion, and the theory of astronomical refraction. We hope that it will soon be put to the full use it deserves. --- Observations for determining the solstices and the obliquity of the ecliptic in 1810 and 1811, by \textsc{Angelo Cesaris}. While \textsc{Oriani} observed the winter solstice of 1810 and the summer solstice of 1811 with the three-foot \textsc{Reichenbach} circle, \textsc{Cesaris} used the eight-foot mural quadrant for the same purpose. As reported here, their results agree within a few tenths of a second. We would only have wished for an explanation of how the quadrant's collimation error was determined, which \textsc{Cesaris} assumes to be $= -1''5$ at both solstices. Nor do these observations indicate any difference in obliquity between the winter and summer solstices, as some astronomers claim to have found. Under the assumption of the same obliquity, \textsc{Cesaris} derives a pole elevation of $45^\circ 28' 0''20$, which agrees very closely with the value found by another method. --- On the degree of convergence of the different series that represent the inequalities in the Moon's longitude, by \textsc{Franz Carlini}. An exceptionally valuable paper! Three expressions for the lunar inequalities are brought together here. The first follows \textsc{Tobias Mayer}'s form, in which \textsc{Bürg}'s tables were calculated. The second, derived from the first by \textsc{Laplace}, is in the form in which \textsc{Laplace}'s theory gave these inequalities, that is, as functions of the Moon's true longitude. The third, calculated by \textsc{Carlini} from the first, is in the form to which \textsc{Lambert} and \textsc{Schulze} had already reduced the equations of \textsc{Mayer} and \textsc{Mason}, that is, as functions of the mean motions alone. We would prefer this last form to all the others, even though the expression converges somewhat more slowly than the first two. It would be desirable for \emph{all} its coefficients to be derived directly from a few thousand observations. Of course, only astronomers who had the preliminary work of \textsc{Bürg} or \textsc{Burckhardt} at their disposal could do this to advantage. \textsc{Carlini} himself tested the expression on 100 positions of the Moon derived from stellar occultations and found good agreement throughout. A noteworthy result is \textsc{Carlini}'s confirmation here of a new equation that \textsc{Burckhardt} had proposed earlier but has left out of his new tables. It first makes its appearance as a long-period inequality in the eccen\SourceLeaf{563}tricity. --- On the effect of changes in temperature on the motion of the pendulum, by \textsc{Carl Brioschi}. The author used a more elaborate method than the investigation required. In essence the result is simply that, in a certain interval of time, a pendulum whose length varies with the temperature makes exactly as many oscillations as it would have made at a constant length corresponding to the mean thermometer reading. --- Finally, some stellar occultations observed in Milan and Rome by \textsc{Carlini} and \textsc{Oriani} in 1808 and 1810.

The supplements to the volume for 1813 begin with the continuation mentioned above of the zenith distances observed with the three-foot \textsc{Reichenbach} circle. --- On the periodic swaying of buildings, by \textsc{Angelo Cesaris}. It has long been known how little one can rely on the position of astronomical instruments that need a fixed mounting to remain unchanged when they are connected to the outer walls of a building, especially at a considerable height above the ground. Temperature and humidity act directly on those walls. It is therefore important, when building a new observatory, to give such instruments a perfectly stable foundation at ground level, entirely independent of the outer walls. The Brera observatory in Milan, otherwise equipped with such excellent instruments, lacks this advantage. The building's movements are very clear and pronounced in the mural quadrant and the transit instrument. If the vertical wire of the transit instrument touches the meridian mark set up at a distance of 3000 toises in the morning before sunrise, the wire subsequently moves farther and farther east of the mark until after noon, when the separation reaches 5 to 6 seconds on clear winter days and 30 seconds on clear sunny days. Other things being equal, the change is smallest or wholly imperceptible when the sky is overcast and the temperature stays nearly constant, and greatest in sunshine with a large change in temperature. --- On the formulas for the parallax and latitude of the Moon, by \textsc{Franz Carlini}, a continuation of the paper in the preceding volume. The parallax formula is transformed, in much the same way as the formula for the inequalities in longitude, into another that depends only on mean motions. \textsc{Carlini} has developed a similar expression for the logarithm of the parallax. He did not likewise transform the expression for latitude, since it would have become far too unwieldy and slow to converge. \textsc{Carlini} confines himself to giving a very simple way of conveniently arranging tables for \textsc{Laplace}'s form (as functions of the true longitude). For several reasons we would still have liked Carlini, for this element too and \emph{with only the first two terms excluded}
\[
18540''25\sin\delta + 12''56\sin 3\delta\text{,}
\]
to have made that transformation. That would have removed the difficulty, and the analogy between the formulas for the inequalities of the Moon and those for the inequalities of the planets would thereby have been complete. The transformation can still easily be carried out. --- Opposition of Saturn in 1811, observed by \textsc{Carl Brioschi}. --- Extract from the meteorological observations at the Milan observatory in 1809, by \textsc{Angelo Cesaris}.

